\documentclass[aps,prb,twocolumn,showpacs]{revtex4-1}
\usepackage{amssymb}
\usepackage{graphicx}
\usepackage{amsmath}

\usepackage{times}
\usepackage{color}
\usepackage{subfigure}
\usepackage{setspace}
\usepackage{epstopdf}
\usepackage{bm}% bold math

%--------------------------------------%
\newcommand{\rcol}{\textcolor{red}}
\newcommand{\beq}{\begin{equation}}
\newcommand{\eeq}{\end{equation}}
\newcommand{\bal}{\begin{aligned}}
\newcommand{\eal}{\end{aligned}}
\newcommand{\bra}{\langle}
\newcommand{\ket}{\rangle}
\newcommand{\vS}{\vec{S}}
\newcommand{\vh}{\vec{h}}
\newcommand{\vQ}{\vec{Q}}



%--------------------------------------%

\begin{document}

\title{Disordered $J-J'$ Heisenberg Antiferromagnet on a Triangular Lattice}
\author{Santanu Dey and Matthias Vojta}
\affiliation{Institut f\"{u}r Theoretische Physik,
Technische Universit\"{a}t Dresden, 01062 Dresden}

\pacs{}
\date{\today}

\begin{abstract}

    The effect of quenched bond disorder in the triangular lattice 
    Heisenberg antiferromagnet has attracted a lot of attention lately.
    Recent experiments on crystals which are believed to realize this model
    typically suggest emergence of a glassy phase when the exchange 
    interactions between the spin degrees of freedom are sufficiently 
    disordered. It is also understood that in the presence of bond disorder
    the $120^{\circ}$ non-collinear chiral order is destroyed in d = 2 lower 
    critical dimension, although, with a correlation length that scales 
    exponentially with the inverse power of the width of the bond disorder. 
    This rules out the possiblity to characterize the system without carefuly 
    eliminating the finite size effects. To this end we resort to two 
    different microscopic modelling of the disordered $J_1-J_2$ Heisenberg 
    hamiltonian, improving on the system size limitations of prior 
    calculations and also mapping out the entire phase diagram of the 
    model which corroborates to the signatures observed in experiments.

\end{abstract}

\maketitle

\section{Introduction}

The heisenberg antiferromagnet in a triangular lattice has long been a 
model of interest as it realizes the simplest case of geometric frustration
in a spin model. The quantum mechanical next nearest neighbour $J_1-J_2$ 
Heisenberg hamiltonian in thought to produce an interesting phase diagram
with a $Z_2$ spin liquid phase interpersed between the non-collinear $120^{\circ}$
spiral and a stripy magnetic order.

\section{The disordered $J-J'$ Heisenberg Hamiltonian on a triangular lattice}

The disordered $J-J'$ antiferromagnetic Heisenberg model in the presence of an
external field can simply be expressed in terms of the bilinears of spin 
operators,
\beq
    \bal
    H &= \sum_{\bra i,j \ket} J_{i,j}\vS_i\cdot\vS_j 
    + \sum_{\bra\bra i,j \ket\ket}\alpha J'_{i,j}\vS_i\cdot\vS_j
    + \sum_i\vh_i\cdot\vS_i 
    \eal
\eeq
where the symbols $\bra i,j\ket$ and $\bra\bra i,j\ket\ket$ represents sum over
nearest and next to nearest neighbours in the triangular lattice. In a clean
system all the couplings will be fixed to a scale and in the absence of the external 
field $\vh$, $\alpha$ becomes the sole tunable parameter that determines the relative 
strength of the next to nearest neighbour (NNN) couplings over the nearest neighbour 
(NN) couplings. However, in a disordered lattice the couplings will be random and this
can be simulated by having $J$ and $J'$ drawn from a random distribution $P(J,\Delta)$ 
with $\Delta$ characterizing the width of the disorder distribution. $\Delta$ is also 
called the disorder strength as the wider the distribution gets the less likely it 
becomes to find a neighbourhood of uniform couplings in the lattice. The parameter 
$\alpha$ now underscores the relative strength between the average NNN couplings and 
the average NN couplings.

The clean zero field limit ($\Delta=0,\vh=0 $) has been studied extensively using 
various microscopic and field theoretic methods. Classically, it can be shown that 
the ground state of the hamiltonian in this limit goes through a phase transition 
from the $120^{\circ}$ non-collinear magnetic order with an ordering wave vector 
$\vQ=(2\pi/3,0)$ to a collinear stripe order with \rcol{$\vQ=(x,x)$}. 

\section{Spinwave excitations with finite disorder}

\subsection{Zero magnetic field}

\subsection{Large magnetic field}

\section{Mean-field theory in real space}

\section{Results}

\section{Conclusions}
\end{document}
